% Tento soubor nahraďte vlastním souborem s obsahem práce.
%=========================================================================
% Autoři: Michal Bidlo, Bohuslav Křena, Jaroslav Dytrych, Petr Veigend a Adam Herout 2019

% Pro kompilaci po částech (viz projekt.tex), nutno odkomentovat a upravit
%\documentclass[../projekt.tex]{subfiles}
%\begin{document}

\chapter{Úvod}
\label{introduction}
\todo{NAPSAT UVOD :(}

\chapter{Počítačové vidění}
\label{computer_vision}
Počítačové vidění je obor umělé inteligence, využívající strojového učení a neuronových sítí, k digitálnímu zpracování obrazu a rozpoznávání vzorů. Jedná se o rozsáhlou disciplínu, která je blízce provázána s počítačovou grafikou a zpracováním obrazu, a jejíž výstupem je porozumnění obrazu \cite{cv_Wiley2018}.

V určitém slova smyslu je počítačové vidění postavené na stejném principu, jako lidské oko. Rozdíly ve schopnosti analýzy obrazu pomocí 'umělého oka' jsou však stále velmi patrné, a to především díky omezené výkonnosti ve srovnání s okem lidským \cite{cv_Wiley2018}. Výhodou lidského zraku je celoživotní kontext, díky kterému se může naučit rozeznávat objekty a jejich zasazení do scény \cite{cv_ibm}.

Tato kapitola se zaměřuje na jednu z mnoha oblastí počítačového vidění, konkrétně na analýzu obrazu, umožňující vyhledávání ve videu.

\section{Vyhledávání ve videu}
\todo{Co to vlastně je vyhledávání ve videu... Proč potřebujeme vyhledávat v záznamech - realtime analýza x potřeba záznamu později, dlouhé záznamy a potřeba specifické události bez znalosti přesného času}
\subsection{Existující systémy}
\todo{...}
\subsection{Problémy spojené s vyhledáváním}
\todo{Faktory ovlivňující vyhledávání - kvalita, překrytí, ...}
\subsection{Etika}
\todo{Zneužití takového systému - oproti realtime monitorování je nebezpečí zneužití mnohem menší}

\chapter{Rozpoznávání objektů ve videu} 
\label{object_detection}
\section{Zpracování videa na snímky}
\todo{ffmpeg vyuziti pro rozkouskovani videa, zmena rozliseni, vybirani vhodneho poctu snimku pro maximalizaci efektivity za cenu zachovani presnosti}
\section{Principy neuronových sítí}
\todo{Mozna ani nedavat???}
\section{You Only Look Once model}
\todo{Popis principu fungovani YOLO modelu a jeho vyuziti souvisejici s vyhledavanim}
\subsection{Ohraničující obdelníky}
\todo{YOLO model vytvarejici bounding boxes}
\subsection{Segmentace}
\todo{YOLO model pro segmentaci - vyuziti pri detekci - cena x presnost}
\section{Transformery}
\todo{Princip moderniho fungovani modelu - nejrozsirenejsi architektura... Zamerit se hlavne na ViT}
\section{Způsoby vyhodnocení}
\todo{Popsa a vysvetlit ruzne zpusoby vyhodnocovani modelu/systemu pro vyhledavani - cena, presnost, ...}


\chapter{Sledování objektů}
\label{object_tracking}
\todo{Co to je object tracking, k cemu je napomocny a jake data se daji ziskat navic s jeho vyuzitim, jaky je obecny princip}
\section{ByteTrack}
\todo{Jakozto hlavni algoritmus pouzity - popsat vyhody, nevyhody, princip fungovani}
\section{deepSORT}
\todo{Sekundarni moznost algoritmu - popsat vyhody, nevyhody, princip fungovani}
\section{Srovnání algoritmů}
\todo{Sorvnání obou algoritmu, jaky je vhodnejsi pouzit a proc....}

\chapter{Zpracování celé scény}
\label{scene_embedding}
\todo{Pokrocile zpracovani cele sceny a ne pouze detekce objektu}
\section{X-CLIP}
\todo{Popsat jak funguje X-CLIP model, zminit i CLIP na kterem je postaveny - vector embeddings, ...}
\section{Vektorový prostor}
\todo{Popsat jak se pracuje s embeddings ziskanymi pomoci X-CLIP, jak se ukladaji informace, ...}


\chapter{Zpracování dotazu pro vyhledávání v přirozeném jazyce}
\label{query_parsing}
\todo{Opravdu jen nejake kratke overview zpracovani prirozeneho jazyka pro uceleni konceptu}
\section{SigLIP}
\todo{Popsat vyuziti modelu SigLIP a obecny princip}
\subsection{Tokenizace a vyhledávání v indexu}
\todo{Zpusob tokenizace za vyuziti SigLIPu a sestaveni dotazu pro vyhledavani cilovych objektu v databazi}
\subsection{X-CLIP embedding dotazu pro vyhledávání}
\todo{Jak se vytvori z query embedding pro nasledne vyhledavani ve vektorovem prostoru}
\subsubsection{Kosínová podobnost (Cosine similarity)}
\todo{Hledani na zaklade podobnosti - similarity search .... cosine similarity}


\chapter{Architektura a implementace}
\label{implementation}
\section{Návrh systému}
\todo{Kratky popis systemu, zpusob navrhu, kratky popis PROC prave takto.}
\section{Architektura systému}
\todo{Navrh systemu - diagramy. 3-urovnove zpracovani podle potreby}
\section{Implementace systému}
\todo{V cem implementovano, co pouzito}
\subsection{1. úroveň zpracování}
\todo{Prvni uroven pro detekci objektu + barev - YOLO}
\subsection{2. úroveň zpracování}
\todo{Druha uroven pro detekci objektu, zpresneni barev, smer pohybu - YOLO + ByteTrack/deepSORT}
\subsection{3. úroveň zpracování}
\todo{Treti uroven pro zpracovani cele sceny za pomoci embeddings - X-CLIP}

\chapter{Testování a vyhodnocení}
\label{tests}
\todo{V teto casti budou popsany pouzite zpusoby testovani, pouzite datove sady, vysledky testovani, ...}
\section{Testování}
\todo{Pouzite metody testovani a jak}
\section{Vyhodnocení}
\todo{Vysledky testovani systemu}


\chapter{Závěr}
\label{zaver}


%===============================================================================

% Pro kompilaci po částech (viz projekt.tex) nutno odkomentovat
%\end{document}

